\section{Testing Strategy}

\subsection{Overview}
The testing strategy follows a strict Test-Driven Development (TDD) approach, where each transformation module is thoroughly tested before proceeding to the next stage. This ensures data integrity and correct calculations throughout the pipeline.

\subsection{Testing Philosophy}
\begin{itemize}
    \item \textbf{Test-First Development:} Tests are written before implementing the actual functionality
    \item \textbf{Comprehensive Coverage:} Each transformation stage has its own test suite
    \item \textbf{Synthetic Data:} Carefully crafted test data to verify edge cases
    \item \textbf{Dual Implementation Testing:} Both Pandas and DuckDB implementations are tested
\end{itemize}

\subsection{Test Data Generation}
Test data is created with specific characteristics:
\begin{itemize}
    \item Small but representative datasets
    \item Known input-output relationships
    \item Edge cases and special scenarios
    \item Consistent date ranges and user patterns
\end{itemize}

\subsection{Test Suites}

\subsubsection{Meal Data Combination Tests}
\texttt{test\_1\_combine\_meal\_data.py}:
\begin{itemize}
    \item Tests meal record combination logic
    \item Verifies correct joining of user and meal data
    \item Validates nutritional calculations
    \item Ensures proper handling of missing data
\end{itemize}

\begin{lstlisting}[style=pythonstyle, caption=Meal Combination Test Example]
def test_meal_combination(self):
    # Test data with known relationships
    test_meals = pd.DataFrame({
        'meal_id': [1, 2],
        'user_id': [1, 1],
        'date': ['2024-04-01', '2024-04-01']
    })
    # Verify combined results
    assert_frame_equal(result, expected)
\end{lstlisting}

\subsubsection{Data Aggregation Tests}
\texttt{test\_3\_group\_data.py}:
\begin{itemize}
    \item Tests daily and user-based aggregations
    \item Verifies group-by operations
    \item Validates statistical calculations
    \item Tests both Pandas and DuckDB implementations
\end{itemize}

\begin{lstlisting}[style=pythonstyle, caption=Aggregation Test Example]
class TestPandasAggregation(TestDataSetup):
    def test_daily_totals(self):
        aggregator = PandasAggregation(self.test_dataframe)
        result = aggregator.compute_daily_totals()
        # Verify daily calculations
\end{lstlisting}

\subsubsection{Window Function Tests}
Three test suites for different window analyses:

\paragraph{Daily Windows} \texttt{test\_4\_1\_window\_function\_daily.py}:
\begin{itemize}
    \item Tests daily rolling calculations
    \item Verifies date-based windows
    \item Validates statistical aggregations
\end{itemize}

\paragraph{User Windows} \texttt{test\_4\_2\_window\_function\_user.py}:
\begin{itemize}
    \item Tests user-specific patterns
    \item Verifies rolling averages by user
    \item Validates temporal sequences
\end{itemize}

\paragraph{Food Type Windows} \texttt{test\_4\_3\_window\_function\_type\_food.py}:
\begin{itemize}
    \item Tests food category analysis
    \item Verifies nutritional proportions
    \item Validates category groupings
\end{itemize}

\subsubsection{Percentage Change Tests}
Two test suites for change analysis:

\paragraph{Daily Changes} \texttt{test\_5\_1\_percentage\_change\_daily.py}:
\begin{itemize}
    \item Tests weekday-based changes
    \item Verifies 4-week rolling averages
    \item Validates percentage calculations
\end{itemize}

\begin{lstlisting}[style=pythonstyle, caption=Percentage Change Test]
def test_percentage_changes(self):
    # Test data spanning 4 weeks
    dates = pd.date_range('2024-03-01', '2024-03-28')
    # Verify percentage calculations
    assert_frame_equal(result, expected)
\end{lstlisting}

\paragraph{User Changes} \texttt{test\_5\_2\_percentage\_change\_user.py}:
\begin{itemize}
    \item Tests user-specific changes
    \item Verifies 4-day rolling averages
    \item Validates individual patterns
\end{itemize}

\subsection{Testing Best Practices}

\subsubsection{Test Structure}
Each test suite follows a consistent pattern:
\begin{enumerate}
    \item Setup test data in \texttt{setUpClass}
    \item Define expected results
    \item Execute transformations
    \item Assert equality with expected results
    \item Cleanup resources
\end{enumerate}

\subsubsection{Assertion Strategies}
Different types of assertions are used:
\begin{itemize}
    \item \texttt{assert\_frame\_equal}: For DataFrame comparisons
    \item \texttt{assertEqual}: For scalar values
    \item \texttt{assertAlmostEqual}: For floating-point comparisons
    \item \texttt{assertRaises}: For error handling
\end{itemize}

\subsection{Test Data Examples}
Example of synthetic test data:

\begin{table}[h]
    \centering
    \begin{tabular}{|l|l|l|l|l|}
        \hline
        \textbf{meal\_id} & \textbf{user\_id} & \textbf{date} & \textbf{calories} & \textbf{protein} \\
        \hline
        1 & 1 & 2024-04-01 & 500 & 20 \\
        2 & 1 & 2024-04-01 & 700 & 30 \\
        3 & 2 & 2024-04-01 & 600 & 25 \\
        \hline
    \end{tabular}
    \caption{Example Test Dataset}
\end{table}

\subsection{Quality Assurance Process}
Testing workflow:
\begin{enumerate}
    \item Write tests for new functionality
    \item Implement the feature
    \item Run all tests to ensure no regressions
    \item Verify edge cases and error handling
    \item Document test cases and expected behavior
\end{enumerate}

This comprehensive testing strategy ensures the reliability and accuracy of the nutritional data pipeline, with each transformation stage thoroughly validated before proceeding to the next phase.
